% Respuesta cualitativa de una celda de carga subamortiguada al impacto de
% una masa: pico, oscilación y asentamiento. Sin escala numérica.
\begin{tikzpicture}[
    >=Latex,
    font=\small,
    eje/.style={-{Latex[width=5pt,length=5pt]}, thick},
    nota/.style={font=\footnotesize, align=center}
]
\draw[eje] (0,0) -- (10.5,0) node[below left] {Tiempo};
\draw[eje] (0,0) -- (0,4.2) node[above right] {Lectura de la celda};

% Masa real
\draw[densely dashed, gray] (0,2.2) -- (10.2,2.2);
\node[nota, anchor=east] at (0,2.2) {Masa real};

% Señal: cero, impacto y oscilación amortiguada
\draw[very thick] (0,0) -- (1.0,0) -- (1.15,3.7)
    -- plot[domain=1.15:10, samples=200, smooth]
    (\x, {2.2 + 1.5*exp(-0.55*(\x-1.15))*cos(deg(4.2*(\x-1.15)))});

% Banda de tolerancia
\draw[gray] (0,2.05) -- (10.2,2.05);
\draw[gray] (0,2.35) -- (10.2,2.35);
\node[nota, anchor=west] at (10.25,2.2) {Tolerancia};

% Rótulos
\node[nota, anchor=west] at (1.35,3.6) {Impacto del material};
\node[nota] at (3.3,0.6) {Oscilación};
\draw[gray] (5.4,0) -- (5.4,2.05);
\node[nota, anchor=north] at (5.4,-0.05) {Asentamiento};
\end{tikzpicture}
